\section{Introduction}\label{sec:intro}

Unstructured kinetic models remain the working description of fed-batch bioprocesses. They lump the culture into a few extracellular concentrations---biomass, substrates, products---and close the mass balances with algebraic rate laws for growth, uptake and product formation \citep{Villadsen2011}. Such models are small enough to be calibrated from routine offline measurements, interpretable enough to be trusted for scale-up and process design, and cheap enough to be embedded in optimisation and model-predictive control \citep{Almquist2014,Narayanan2020}. Their usefulness, however, rests on two decisions that are still made largely by hand: which rate law governs the culture, and what its parameters are.

Both decisions are usually resolved by fitting. Each candidate law---Monod saturation \citep{Monod1949}, Contois-type biomass dependence \citep{Contois1959}, Haldane substrate inhibition \citep{Andrews1968}, product inhibition, cell death and their combinations---is calibrated by nonlinear least squares and the candidates are compared through an information criterion that penalises additional parameters \citep{Akaike1974,Schwarz1978,Burnham2002}. The procedure is statistically sound but practically fragile. Every fit is a non-convex problem whose outcome depends on the starting point; the cost grows with the number of candidates; and individual constants are often only weakly identifiable from a single trajectory \citep{Raue2009,Saa2017}. Two further limitations receive less attention. First, a selection procedure can only return the best of the candidates it is given: if the law that governs the culture is not in the library, the procedure still returns an answer, and nothing in the ranking itself signals that the answer is wrong. Second, a single fed-batch run is often uninformative about \emph{why} growth slows. Biomass rises while the limiting substrate falls, so a law in which the specific rate decreases with biomass and a law in which it changes with substrate can reproduce the same run. Model-based design of experiments addresses this second problem by choosing experiments that separate rival models \citep{HunterReiner1965,BoxHill1967,BuzziFerraris1983,Franceschini2008}, but it requires the rival models to be written down in advance.

Machine learning offers a complementary route. Hybrid, or semi-parametric, models keep the known mass balances and learn the unknown rates with a neural network \citep{Psichogios1992,Thompson1994,Oliveira2004,vonStosch2014,Sansana2021,Bradley2022}; neural ordinary differential equations (neural ODEs) make such models trainable end to end by differentiating through the integrator \citep{Chen2018,Kidger2021}; and physics-informed neural networks embed the governing equations in the training loss \citep{Raissi2019,Bangi2022}. These approaches learn \emph{that} a rate behaves in a certain way, but they do not by themselves say \emph{which} mechanism it is, and a neural network is not the interpretable, extrapolable rate law that process engineers need. Equation-discovery methods close part of this gap. Sparse regression over a dictionary of candidate terms recovers governing equations from data \citep{Brunton2016,Champion2019}, including the rational rate laws typical of biochemical kinetics when the problem is written in implicit form \citep{Mangan2016,Kaheman2020}, and related ideas have been applied to reaction networks \citep{Hoffmann2019,JiDeng2021}; genetic-programming symbolic regression searches over free-form expressions \citep{Schmidt2009,Cranmer2023}. \citet{Rackauckas2020} showed that the unknown term learned by a universal differential equation can subsequently be converted into a closed-form expression by sparse regression. In a fed-batch bioprocess, however, measurements are sparse, noisy and offline, numerical derivatives are unreliable, the admissible rate laws must satisfy physical constraints (non-negativity, vanishing without substrate), and any proposed law must compete on equal statistical terms with the established mechanisms rather than replace them.

In this work we combine these elements into a single workflow for kinetic mechanism identification that \emph{learns the kinetics first and names them afterwards}. A hybrid neural ODE, whose specific rates are a small neural network constrained to be non-negative and to vanish in the absence of substrate, is fitted directly to the measured trajectories, without numerical differentiation and without assuming a mechanism. The learned rates are then used in three ways: (i) to place every candidate mechanism of a library at a data-informed starting point, after which each candidate is refined against the measurements and the candidates are ranked by the Bayesian information criterion; (ii) to propose new rate laws through a symbolic regression whose candidate laws are fitted in rate space and restricted to physically admissible forms, including non-rational ones; and (iii) as a reference when a goodness-of-fit test shows that no candidate explains the data, in which case every admissible law is screened against the measurements and a law that passes both a model-selection and an adequacy criterion is added to the library for later experiments. Throughout, several runs with different inocula, initial substrate and feeds can be fitted jointly with shared kinetics, which decorrelates biomass and substrate and removes the ambiguity of single runs.

The workflow is evaluated on two case studies. The first is a single-substrate fed-batch system with a library of eleven mechanisms, in which the generating mechanism is known and a twelfth, deliberately hidden law is absent from the library; it is used to quantify, over repeated noise realisations, the accuracy of selection from one run and from three jointly fitted runs, the false-discovery rate of the symbolic regression, the discovery of the hidden law and the effect of adding it to the library. The second is a dual-substrate xylitol bioconversion by \emph{Saccharomyces cerevisiae} with a library of eleven mechanisms, studied first on virtual runs that reproduce the feeds, volumes and sampling schedules of real fermentations and then on four real fed-batch runs \citep{Toivari2025}, individually and jointly, with the measurement noise treated as unknown.

The contributions of this paper are: (1) a hybrid neural-ODE formulation of fed-batch kinetics whose rate parameterisation is mechanism-agnostic but physically admissible, trained on sparse offline data; (2) a selection procedure in which learned rates initialise every candidate mechanism and selection is decided by refined fits and information criteria, with a profile-likelihood variant for unknown noise levels; (3) a physically constrained symbolic regression that proposes rational and non-rational rate laws and ranks them on the same footing as the library, together with an adequacy-triggered screening and acceptance rule that grows the library; and (4) a systematic evaluation on synthetic data with known ground truth and on four real fermentations, showing when single runs are ambiguous, how jointly fitted runs resolve the ambiguity, and why the shared kinetics that joint fitting assumes must be tested on real data.

The remainder of the paper is organised as follows. Section~\ref{sec:background} reviews kinetic model discrimination, hybrid models and equation discovery. Section~\ref{sec:methods} presents the workflow. Sections~\ref{sec:case1} and \ref{sec:case2} report the two case studies, Section~\ref{sec:discussion} discusses the results and their limitations, and Section~\ref{sec:conclusions} concludes.
